% Motivation and learning outcomes.
% The outcomes are the outcomes of Day 10 in the syllabus.
\begin{frame}{Why this day matters}
  \begin{itemize}
    \item Days 1 to 9 used models with thousands or millions of parameters. A
          small language model has billions of parameters. Its weights fill a
          large part of the RAM of the Raspberry Pi 5.
    \item A language model writes one token after the other. The speed of a
          chat is the number of tokens in each second, and the memory
          bandwidth limits it more than the arithmetic.
    \item A local model keeps the data on the device and works with no
          network. It can also give a wrong answer with confidence, and it
          uses much power and heat.
    \item In the lab you run two models of different size with Ollama, measure
          them with the method of Day 9, call a model from Python with
          structured output, and build one small application.
  \end{itemize}
\end{frame}

\begin{frame}{Learning outcomes}
  After this day you can:
  \begin{itemize}
    \item Estimate the RAM that a small language model needs from its
          parameters, its precision, and its context length.
    \item Run small language models on the Raspberry Pi and measure tokens per
          second.
    \item Call a small language model from Python with structured output and
          function calling.
    \item Explain retrieval-augmented generation and vision-language models at
          the edge.
    \item State when a small language model is the wrong tool.
  \end{itemize}
\end{frame}
